% !TeX root = ../main.tex
\chapter{JADWAL PENELITIAN}

\section{Rencana Kerja}

Rencana kerja disusun untuk delapan bulan dan dapat disesuaikan setelah hasil pilot, ketersediaan model, serta jadwal akademik ditetapkan. Blok berwarna menunjukkan periode aktivitas utama.

\begin{table}[htbp]
\centering
\small
\renewcommand{\arraystretch}{1.25}
\newcommand{\taskbar}{\cellcolor{blue!30}}
\resizebox{\textwidth}{!}{%
\begin{tabular}{|p{5.0cm}|*{8}{c|}}
\hline
\textbf{Kegiatan} & \textbf{Sep. 2026} & \textbf{Okt. 2026} & \textbf{Nov. 2026} & \textbf{Des. 2026} & \textbf{Jan. 2027} & \textbf{Feb. 2027} & \textbf{Mar. 2027} & \textbf{Apr. 2027} \\
\hline
Studi pustaka dan finalisasi taxonomy & \taskbar & \taskbar & & & & & & \\
\hline
Pemilihan benchmark dan task manifest & & \taskbar & \taskbar & & & & & \\
\hline
Konstruksi treatment dan treatment checker & & \taskbar & \taskbar & \taskbar & & & & \\
\hline
Implementasi loader validation dan logger & & & \taskbar & \taskbar & & & & \\
\hline
Smoke test dan pilot berpasangan & & & \taskbar & \taskbar & & & & \\
\hline
Eksperimen utama dan pengarsipan artefak & & & & \taskbar & \taskbar & & & \\
\hline
Analisis outcome, cost, dan failure mode & & & & & \taskbar & \taskbar & & \\
\hline
Penulisan hasil dan pembahasan & & & & & & \taskbar & \taskbar & \\
\hline
Revisi, validasi reproduksi, dan finalisasi & & & & & & & \taskbar & \taskbar \\
\hline
\end{tabular}%
}
\end{table}

\section{Keluaran Tiap Tahap}

\begin{table}[htbp]
\centering
\small
\renewcommand{\arraystretch}{1.25}
\caption{Keluaran yang ditargetkan}
\label{tab:deliverables}
\begin{tabularx}{\textwidth}{p{4.0cm}Y}
\toprule
\textbf{Tahap} & \textbf{Keluaran} \\
\midrule
Studi pustaka & Taxonomy enam smell, definisi operasional tiga treatment, research question, hipotesis, dan gap. \\
Persiapan benchmark & Task manifest, pinned commit, dependency, command validasi, acceptance criterion, dan lisensi. \\
Konstruksi treatment & File Clean, tiga varian smell, varian Remediated, diff, token count, dan treatment checker. \\
Pilot & Laporan loader validation, reproducibility, variasi outcome, estimasi biaya, dan keputusan go/no-go. \\
Eksperimen utama & Dataset run, patch, trajectory, raw output, validation result, dan resource log. \\
Analisis & Tabel outcome, effect size, interval ketidakpastian, paired comparison, dan failure taxonomy. \\
Finalisasi & Draft proposal atau skripsi, artefak reproducible yang aman dipublikasikan, serta dokumentasi batas klaim. \\
\bottomrule
\end{tabularx}
\end{table}
